\section{Problem Statement}
\label{sec:problem}

We now state the problem precisely. A specification defines a set of
operations \Ops. A tester sends a sequence of at most \Budget requests
$r_1, \dots, r_{\Budget}$, each for one operation $\mathit{op}(r_i) \in \Ops$,
and observes a status code $\mathit{status}(r_i)$. The tester chooses each
request after seeing the responses to the earlier ones.

\paragraph{Reached operations.} An operation is \emph{reached} when at least
one request for it returned a 2XX status:
\begin{equation}
  \Reached = \{\, o \in \Ops \mid \exists i:\ \mathit{op}(r_i) = o \ \wedge\
  200 \le \mathit{status}(r_i) < 300 \,\}.
  \label{eq:reached}
\end{equation}
$|\Reached|$ is the operation coverage of REST League.

\paragraph{Unique faults.} A function \sig maps a 5XX response to a
\emph{signature}: the operation, the status code and the error message with
the parts that change from request to request removed (numbers, ids, and
values the API echoes back in quotes). The unique faults are
\begin{equation}
  \Faults = \{\, \sig(r_i) \mid 500 \le \mathit{status}(r_i) < 600 \,\}.
  \label{eq:faults}
\end{equation}

\paragraph{Goal.} Given the specification and the budget \Budget, produce
requests that make $|\Reached|$ as large as possible, and reach it with as
few requests as possible, so that the rest of the budget can go to error
requests that grow $|\Faults|$. This report studies the input side of this
goal: which values to put into each request, and in which order to try the
operations. Choosing what to send next during a run is the engine's job and
is outside this report.

\paragraph{Why is this hard?} A schema states only part of what the API
checks. The API also checks rules the schema does not state: a date string
must parse, an id must name an existing resource, and a resource must exist
before it is read. A request that follows the schema can still be rejected,
and a request that breaks it can still be accepted.

\paragraph{Why not just send the examples in the specification?} Because
most specifications have none. In our set of \NSpecs specifications,
\NSpecsWithoutExamples contain no \texttt{example} value at all
(Table~\ref{tab:parser} lists the specifications). Even where examples exist,
they cover only some parameters, and an example id such as \texttt{1} is
right only while resource~1 exists.
